\section{Threat Model}
\label{sec:threat}

Because \defname{} is an environment, its threat model is \emph{parameterized and
configurable}, not fixed. A study selects a point in the adversary parameter space, and
the harness realizes it. \Cref{fig:threat} depicts the resulting honest, adversarial,
and orchestrator/tool/memory vectors and how each enters the system.

\subsection{Adversary configuration}
An adversary is specified by a configuration
\begin{equation}
\label{eq:adv-config}
\config_{\mathrm{adv}} \;=\; \bigl(\,\betafrac,\ \mathsf{coll},\ \mathsf{adapt},\
\mathsf{orch},\ \mathsf{tool},\ \mathsf{syb},\ \mathsf{mem}\,\bigr),
\end{equation}
whose components the harness exposes as independent knobs:
\begin{itemize}[leftmargin=1.4em]
\item \textbf{Adversary fraction} $\betafrac=|\Advset|/N\in[0,1]$: the fraction of agents
placed under adversarial control, $\Advset\subseteq\Agents$.
\item \textbf{Collusion} $\mathsf{coll}\in\{\text{independent},\text{colluding}\}$:
whether adversarial agents optimize a shared objective and may coordinate messages,
timing, and target selection, or act independently.
\item \textbf{Adaptivity} $\mathsf{adapt}\in\{\text{static},\text{adaptive}\}$: whether
the red strategy is fixed in advance or may react to the blue defense across \corb{}
rounds (\Cref{sec:algorithm}).
\item \textbf{Orchestrator compromise} $\mathsf{orch}\in\{\text{trusted},\text{compromised}\}$:
a single bit flipping the planner from honest to adversarial. A compromised orchestrator
can mis-route, mislabel, or fabricate subtasks and messages.
\item \textbf{Malicious tools} $\mathsf{tool}$: a subset of tools whose (emulated)
behaviour is adversarial, realized by the adversarial emulator $\Emu^{\!\dagger}$
(e.g., returning attacker-controlled content or exfiltration bait).
\item \textbf{Sybil identities} $\mathsf{syb}$: the number of additional identities an
adversary may instantiate to gain disproportionate influence over selection, voting, or
memory writes (cf.\ the classical Sybil attack~\cite{douceur2002sybil}).
\item \textbf{Poisoned memory} $\mathsf{mem}$: adversarial entries written into the
shared-memory store, in the spirit of memory-poisoning attacks on single
agents~\cite{chen2024agentpoison} but now readable by many agents.
\end{itemize}

\subsection{Attacker objective, knowledge, and capabilities}
\paragraph{Objective.} The adversary seeks \emph{system-level compromise}: causing an
unsafe action, exfiltrating protected information, or sabotaging coordination so the
collective fails a forced-coordination task. Formally, the adversary maximizes
$\Prob[\Predsec(\Trace)=1]$ (attack success) and, for sabotage objectives, minimizes
$\Prob[\Predutil(\Trace)=1]$ (benign success), where $\Predsec,\Predutil$ are the fixed
predicates of \Cref{sec:problem}.
\paragraph{Knowledge.} \defname{} supports a spectrum. At minimum the adversary knows its
own role, tools, and local context; at maximum (a strong-adversary study) it knows the
topology $\Gtop$, the honest agents' prompts, and the active blue defense. Knowledge is a
configurable dimension so that a defense is not credited for security that evaporates
under a better-informed attacker.
\paragraph{Capabilities.} An adversarial agent may emit arbitrary messages and tool calls
within its granted permissions, write to memory it can write to, and---if colluding---share
state with other adversarial agents out of band. A compromised orchestrator additionally
controls routing and task assignment. Adversaries \emph{cannot} read observations they are
not permitted to see, forge messages on edges they do not control (the harness attributes
messages to their senders), or alter the recorded trace $\Trace$; these are TCB
guarantees.

\subsection{Static versus adaptive, independent versus colluding}
The two axes that distinguish the multi-agent regime are collusion and adaptivity. An
\emph{independent} adversary of fraction $\betafrac$ behaves like $\betafrac N$ separately
acting bad agents; a \emph{colluding} adversary of the same fraction coordinates seed
placement, message content, and timing toward the shared objective. A \emph{static}
adversary commits to a strategy; an \emph{adaptive} adversary co-evolves against the blue
policy. \Cref{sec:theory} shows that in the environment's propagation model, collusion
yields a strict advantage on star and tree topologies; \Cref{sec:experiments} tests
whether that advantage manifests as measured defense degradation for real defenses
(a \emph{hypothesis}).

\subsection{Trust assumptions and scope}
Trusted: the deterministic checker $\Checker$ and the harness (\Cref{sec:tcb}). Untrusted
under the appropriate knob: any agent (including the orchestrator), any tool, any memory
entry, and any message content. Out of scope: real-world tool side effects, physical
actuation, and model-weight or training-time attacks on the backbones. We emphasize the
stance of the whole paper: a configuration that yields low $\ASRs$ or high $\NRPs$
demonstrates \emph{measured} risk under the stated knobs; it is never read as a proof of
security.
